% language=uk

\usemodule[art-01,abr-02]

\startdocument
  [title=What is \CONTEXT,
   author=Hans Hagen,
   affiliation=PRAGMA ADE,
   location=Hasselt NL,
   date=\currentdate,
   website=www.pragma-ade.nl,
   wiki=www.contextgarden.net,
   metadata:title=What is ConTeXt,
   metadata:author=Hans Hagen]

\startMPpage

    StartPage ;
        fill Page enlarged 2mm withcolor .4blue ;
        draw anchored(textext("\bf ?") ysized (bbheight Page -2cm),center Page) withcolor .6red ;
        draw anchored.lrt(textext("\bf context") xsized 8cm,lrcorner Page shifted (-1cm,10cm)) withcolor .8white ;
    StopPage ;

\stopMPpage

\setuphead[subject][color=darkblue]

\startsubject[title=Rationale]

Occasionally I run into a description of \CONTEXT\ that contains
observations that are somewhat off. It therefore makes sense to
provide some insight in why this macro package looks the way it
looks. What started out as a \TEX\ only system evolved via adding
\METAPOST\ to the current hybrid system that also uses and provides
\LUA. However, the original goals and principles are still valid.

\stopsubject

\startsubject[title=The system]

The \TEX\ macro language and interpreter are about automated
typesetting. A collection of predefined macros is called a macro
package and \CONTEXT\ is such a package. The program that
interprets the macros and converts input into typeset results is
called the engine. An example of an engine is \LUATEX. The graphic
companion of \TEX\ is \METAFONT, or nowadays \METAPOST\ and
\CONTEXT\ integrates support for that graphical language and
rendering. As \CONTEXT\ comes with a bunch of tools that manage
the process it is a sort of typographic ecosystem. The input can
be a document encoded in structured \TEX\ code or it can be \XML.
In fact anything that makes sense can be processed if only because
we can use \LUA\ to convert it runtime. The output is in most
cases \PDF, but \XML\ (or \XHTML\ with CSS) is also an option. In
addition to static printable documents you can produce highly
interactive documents for screen and reader.

The \TEX\ system, and therefore \CONTEXT, are quite capable to
produce high quality output, can handle many scripts and
languages, are able to typeset mathematics and can keep up pretty
well with developments, thanks to the open character of the whole
system. There is an active community that takes care of the
software distributions and develops the engines as well as
additional resources like fonts. User groups play an important
role in this. The \CONTEXT\ system is supported by active mailing
lists and a wiki.

\stopsubject

\startsubject[title=The syntax]

The syntax is rather straightforward and uses commands that start
with a \tex{}, the curly braces \type {{}} for delimiting variable
content and the square brackets \type {[]} for (optional)
directives and settings. A good example is one of the oldest
commands:

\startTEX
\startitemize[packed]
\item some text
\item some more text
\stopitemize
\stopTEX

Before we ran into \TEX\ we already had some programs that dealt
with simple formatting and start|/|stop constructs were used
there. The not delimited \type {\item} command was just taken from
what we saw in other \TEX\ code. In retrospect we should have used
\type {\startitem} and \type {\stopitem} commands right from the
start but it's only with \MKIV\ that we start promoting this more
strongly. Actually there are only a few commands that are not
delimited and this is one of them. So, this is today's fashion:

\startTEX
\startitemize[packed]
\startitem some text \stopitem
\startitem some more text \stopitem
\stopitemize
\stopTEX

One reason for getting rid of the few non delimited cases is that
is't more cleaner in the perspective of hooking in code as well as
exporting to for instance \XML.

The (optional) argument between square brackets use those brackets
because for a beginning macro writer (which I definitely was when
I started with \CONTEXT) it's not that hard to implement. It also
looks nice on screen, especially when you use syntax highlighting.
When \CONTEXT\ showed up we still used our own editor and it used
a color scheme that can still be found in the pretty printer code.
So, to a large extent, the syntax has been determined by how well
it could be visualized in an editor. We set up commands using:

\startTEX
\setupsomething
  [whatever]
  [key-1=value-1,
   key-2=value-2]
\stopTEX

Spaces after a comma are ignored contrary to spaces before and
after an equal sign. It takes more (and slower) code to do that
and picking up keys and values is already slow enough, especially
when we started doing it.

I'm writing this in the \SCITE\ editor (2.29) using updated syntax
highlighting code that was added to the distribution around the
\CONTEXT\ conference in 2011. Of course the \TEX\ code gets
highlighted, but so does \METAPOST\ and \LUA\ code. Runtime spell
checking is included. Developments like that have some impact on
the user interface and the way (low level) code is written.
Visualization of code and source has been and always will be an
important aspect of how \CONTEXT\ evolves.

\stopsubject

\startsubject[title=The multilingual interface]

A second important aspect of \CONTEXT\ is that there is just one
version. I simply could not come up with a good reason to use
different names for different engines. After all, most of the
(older) core functionality is shared. It is always a surprise to
run into descriptions of \CONTEXT\ being dependent on \PDFTEX,
\PERL, or \RUBY. In \CONTEXT\ there has always been an abstract
layer between the core code and the backend and support and in \MKII\
relatively little amount of backend specific code is loaded. In \MKIV\
we have several backends running at the same time: \PDF\ and \XML\
export can run in parallel.

Being closely involved in the development of \PDFTEX\ naturally
means that I used that engine but originally we used \DVIPSONE\
and \CONTEXT\ \MKII\ supports several backends. The dependency on
\PERL\ and later \RUBY\ was only true for the script that manages
a run (originally the index sorter was written in \MODULA2). One
could run unmanaged but it's much more convenient for users not to
worry about how many runs are needed. It also gave us the
opportunity to provide command line arguments. And of course index
sorting has to happen somewhere out of \TEX\ itself so why not
combine these tasks in one script. When we moved on to \LUA\ it
was natural to stick to only \LUA\ and as \LUATEX\ is also a \LUA\
interpreter, there is no extra dependency in \CONTEXT\ \MKIV.

We have only one version, but on your system you might find
more than one \type {cont-*} file. The reason for this is that
there are different user interfaces. We started with the Dutch
interface and later added an English and German one. Of course
some more followed.

When we started using \LUATEX, it quickly became clear that we had
to split the code base and so we did. In fact one can now indeed
claim a dependency: on \LUATEX. The fact that \MKII\ is frozen is
a clear signal that we see no future for the other engines in the
\CONTEXT\ community and their development is stalled anyway. The
same is true for related technologies like fonts. Why stick to
obsolete font formats when we can move on. Fortunately users are
quite willing to move on with us.

The dependency on \LUA\ is not so much a dependency as well as
progress, especially if we keep in mind that the development of
\LUATEX\ has been driven by the \CONTEXT\ people. It's way more
fun to use the \TEX, \LUA\ and \METAPOST\ languages for what they
are best suited for than to try to squeeze all functionality out
the macro language only. It also suits the original design goals
that Don Knuth gave \TEX: to take the code and adapt it to ones
needs.

An interesting side effect of the multilingual interface is that
it made the low level key/value handling more efficient. Given
that we nowadays have faster machines in \MKIV\ some of the gain
in speed has been sacrified to a more neat but slower inheritance
subsystem. Probably no user will notice this anyway.

\stopsubject

\startsubject[title=Being monolitic]

When you run \CONTEXT\ you don't need to load (or setup) extra
styling code. You will always get some output. For those who come
from other macro packages this is somewhat hard to grasp and
sometimes seen as a disadvantage.

However, a fact is that in practice one can easily modulate on
the defaults.

\startTEX
\setupbodyfont[dejavu]

\setuplayout[width=middle,height=middle]

\setupwhitespace[big]

\setuphead[chapter][style=\bfc]
\setuphead[section][style=\bfb]
\stopTEX

Is the above really much more work than loading a style that
defines the font and another one that sets up the spacing and
styles? Also, a user can put these commands in a file and load
that one. Changing the look and feel this is way more convenient
than loading some default and try to overload unwanted settings
(especially if that style changes). It also gives the user an idea
that there can be a personal touch to the document. Of course the
user can just stick to the defaults.

Any observation that users are supposed to know plain \TEX\ or do
some coding is just wrong and probably come from experiences with
other macro packages. On the other hand it might help the user to
know a bit about the project structure, separating structure and
layout and limiting coding. Much in \CONTEXT\ relates to structure
and the actual rendering is an independent issue. Of course a user
can still do things similar to plain \TEX, so buying a copy of The
\TEX\ Book is no waste: you can use most tricks mentioned there
in \CONTEXT\ and there is a lot of information about fine|-|tuning
math typesetting. It also does not hurt to know a bit about where
we come from.

\stopsubject

\startsubject[title=Speed]

Indeed \CONTEXT\ is not a fast runner, but it's not that slow
either. In some cases a slow terminal is the culprit (as \TEX\
does no buffering), and in other cases the user just asks for
something that needs processing time. Especially decorating the
page will increase the runtime. Of course delegating some action
to \LUA\ costs time, but we gain back functionality that otherwise
would not be possible or take much more runtime. The startup time
of \MKIV\ is much shorter than \MKII\ which is partly due to more
efficient file searching so in practice \MKIV\ runtime is quite
acceptable, espically if we consider that we load larger fonts and
operate in a \UNICODE\ universum. Also, hyphenation patterns are
loaded only when needed and, when used, \METAPOST\ processing
happens instantaneously.

\stopsubject

\startsubject[title=Development]

Indeed most development is done by a few people, but how bad is
that? If we look at the larger picture, there is a whole
infrastructure in place: wiki, stand|-|alone distribution, mailing
lists, conferences, and all hat is taken care of by pretty active
crowd. The advantage of a small development team is that
consistency is easier to guard and (hopefully) the distance
between developers and regular users is rather small. Also, input
with regards to subsystems (language labels, font definitions,
etc.) is accepted from whoever comes up with improvements, given
that they are consistent. There are no formal committees and
reasonable demands are often met within reasonable time.

\stopsubject

\startsubject[title=Specific fields]

There is no reason why you shouldn't use \CONTEXT\ if you need
to typeset math. First of all it runs on top of \TEX\ so you will
get your math done one way or the other. Also, math support is
quite complete although the implementation differs from other macro
packages (this might be even more true in \MKIV). We're not bound
by traditions and have some pretty good experts on board that helps
us to move on. And more is coming.

Another area of interest is critical editions. Again some experts
in the field are involved and \CONTEXT\ keeps evolving in this
area, driven by demands from advanced users.

All the usual bells and whistles that \TEX\ engines can offer,
like character protrusion and hz optimization are supported as are
advanced font features. Also, we already had quite some
specialized functionality for detailed typesetting and much
already has been rewritten in \MKIV. Quite some languages are
supported right out of the box. However, keep in mind that this is
all built in so don't start looking for special styles or so (as
the lack of them might suggest that it's not there). Just join the
mailing list.

\stopsubject

\startsubject[title=Kernel code]

Users normally stick to the more high level commands. There are
however quite some support macros. The real low level code is not
be touched by users. Users coming from other packages might need
to get accustomed to leaving them along and define (and tweak)
instances at the higher level. In \MKIV\ most of \CONTEXT\ is also
available at the \LUA\ end and sometimes solving more complex
problems is done easier that way. By now we have a hybrid system.
As part of the rather drastic \MKIV\ cleanup and rewrite we're in
the process of hiding obscure code from the user and might end up
with better low level \API\ documentation.

On top of the kernel code we have some modules. These implement
very specific functionality and more and more get written. There
are for instance modules for typesetting \MATHML, loading fonts
with less commands, integrating external applications, (advanced)
presentations, tracing. They are good examples of how to write
extensions.

\stopsubject

\startsubject[title=Documentation]

There is quite some documentation but it's rather diverse. The
documents written by the authors are often a side effect of some
development. There is a shift towards more general documentation
(and even printed copies are available) but writing can get a
lower priority in a time when quick and dirty answers can be found
on the internet, mailing list or wiki. On the other hand, the user
interface has always been pretty well defined in a formal way) and
once a user knows what some keys do with a command, he or she can
also use that knowledge for other commands. And, as we aim for
upward compatibility, a decade old manual might still apply so
there is no need to render it again just for the sake of updating
a date.

For what it's worth: whenever I have to solve a problem with a
program (or language implementation) I run into cases where I have
to look long to find (non conflicting) information. It just comes
with the problems one wants to solve and \TEX\ (\CONTEXT) is not
different.

\stopsubject

\startsubject[title=Conclusion]

So let's draw some conclusions and limit ourselves to \MKIV. First
of all this is the versions that (new) users are supposed to use.
It also means that they use \LUATEX, which in turn means that
there are no further dependencies. The \CONTEXT\ suite that you
can download from \type {contextgarden.net} is an easy starting
point and \TEX\ Live is also an option. They should just start and
delay styling to when they feel the need. By that time they only
need a handful of commands to get a decent job done. If you can't
find the documentation you need (on paper of on the wiki),
consider participating in writing it.

\setuplines[before=\vfilll,after=,style=\bf,color=darkblue]

\startlines
\documentvariable{author}
\documentvariable{affiliation}
\documentvariable{location}
\documentvariable{date}
\blank
\documentvariable{website}
\documentvariable{wiki}
\stoplines

\stopsubject

\stopdocument
